% Balade: lifelike cursor behavior from state machines and mathematics.
% Pas de reseaux de neurones : generateurs, pilotage, Fitts, minimum jerk.
\documentclass[11pt,a4paper]{article}
\usepackage[utf8]{inputenc}
\usepackage[T1]{fontenc}
\usepackage{lmodern}
\usepackage{amsmath,amssymb}
\usepackage{hyperref}
\usepackage{booktabs}
\title{Balade: Comportement de curseurs vivant par machines a etats et mathematiques}
\author{Contributeurs Balade}
\date{\today}
\begin{document}
\maketitle
\begin{abstract}
Balade est une extension de navigateur qui pose des curseurs animes autonomes
sur les pages web. Ils se baladent, inspectent les elements, tapent du texte
fantome, se saluent et anticipent l'utilisateur. Aucun reseau de neurones :
chaque comportement est une coroutine (un generateur utilise comme machine a
etats) pilotee par des mathematiques classiques : loi de Fitts, profils a
secousse minimale, geometrie des arcs, ressorts amortis, pilotage de Reynolds,
chaines de Markov et scores de style bayesien. Cet article documente les
formules, les substitutions qui les ont rendues rapides, et les bancs d'essai
qui ont tranche chaque substitution.
\end{abstract}

\section{Comportements comme coroutines}
Chaque curseur execute un generateur \texttt{*run(cursor)} qui rend la main
avec le pas de temps et reprend ensuite. Une seule boucle requestAnimationFrame pompe tous les curseurs with a
shared clock, so reads and writes stay batched. La composition utilise
\texttt{yield*}: \texttt{moveTo}, \texttt{wait}, \texttt{click} are reusable
subroutines. L'etat vit dans les variables locales du generateur, pas dans des champs d'etat.

\section{Mathematiques du mouvement}

\subsection{Duree : loi de Fitts}
\begin{equation}
T = \frac{a + b \log_2(1 + D/W)}{v}, \quad \text{clamped to } [T_{\min}, T_{\max}]
\end{equation}
$D$ est la distance, $W$ la largeur de la cible, $v$ un multiplicateur de vitesse.

\subsection{Profil : secousse minimale avec deformation du temps}
\begin{equation}
s(t) = 10t^3 - 15t^4 + 6t^5, \qquad u = s(t^q), \; q \in [0.68, 0.85]
\end{equation}
Chaque curseur garde sa propre deformation $q$ a vie : une signature motrice.

\subsection{Trajectoires : arcs tortue et similitude}
Tracees facon tortue le long d'arcs de cercle a courbure constante, puis projetees par
une similitude (rotation, echelle, translation) sur la cible. Les trajectoires courtes restent
presque droites ; les longues courbent jusqu'a pleine amplitude.

\subsection{Stabilisation : ressort d'Euler semi-implicite}
\begin{equation}
a = -k\,x - c\,v, \qquad v \mathrel{+}= a\,\Delta t, \quad x \mathrel{+}= v\,\Delta t
\end{equation}

\subsection{Essaim : pilotage de Reynolds}
Orbite a rayon respirant ($\pm 18\%$), rosace a 3 lobes et facteur de tourbillon
$\mathrm{swirl}(t) = \tanh(3\cos(0.11\,t))$ qui inverse la direction environ toutes les 29~s. Vitesse voulue moins vitesse actuelle, force plafonnee, integration d'Euler.
L'alignement avec les voisins melange 70/30 avec la moyenne ponderee des vitesses proches.

\section{Substitutions de formules : meme resultat, moins cher}
Chaque substitution ci-dessous a ete mesuree (200\,000 gestures $\times$ 24 steps,
et micro-bancs cibles). Les alternatives rejetees sont documentees aussi.

\subsection{Rotation iterative du cap ($4\times$)}
Tracer un arc demande $(\cos h, \sin h)$ par pas. Au lieu de deux appels transcendants
par pas, le vecteur de cap unitaire avance par multiplication complexe avec
la rotation fixe $(\cos \Delta h, \sin \Delta h)$: 2 appels par
arc au lieu de 48. Exact au bit (derive $10^{-16}$ sur 24 pas). Rejete : une
table de sinus de 256 entrees ($3\times$ \emph{plus lente} que \texttt{cos} :
l'acces memoire coute plus que le FPU) et Taylor ordre 5 (plus lent \emph{et}
faux, erreur $0.5$ pres de $\pm\pi$).

\subsection{Courbes de progression echantillonnees (zero pow par image)}
Comme $q$ est fixe par geste, $u = \mathrm{minjerk}(t^q)$ est echantillonne une
fois par geste en 64 entrees ; chaque image est une lecture indexee avec
interpolation lineaire. Erreur maximale $5 \times 10^{-4}$ sur la position : sous-pixel.

\subsection{Pilotage triple-angle sans atan2 ($3\times$)}
L'essaim demande $\sin(3\theta + \varphi)$ ou $(\cos\theta, \sin\theta)$ sont
deja connus par la normalisation. Le developpement triple-angle evite
\texttt{atan2} completement (34 contre 100~ms par million d'evaluations).
Les oscillations temporelles sont calculees une fois par image et partagees par tous les curseurs ;
les balayages de voisins utilisent le carre des distances avec \texttt{sqrt} reserve aux proches
($2\times$).

\subsection{Prediction : rester simple}
Un pre-filtre en cosinus ($ \cos < \cos(\mathrm{tol}\cdot 2.5)$) a ete mesure et
rejete : \texttt{cos} coute a peu pres autant que \texttt{acos}, et les survivants paient
les deux (48 contre 38~ms pour 500k candidats). \texttt{acos} direct conserve.

\subsection{Deliberement non optimise}
Fitts's \texttt{log2} tourne une fois par geste : des nanosecondes sur un budget de 16~ms,
et la formule reste lisible. $(1-d/m)^2$ ecrit en toutes lettres mesure identique a
\texttt{pow} (45 vs 43~ms per 2M); kept as $k\cdot k$ pour la lisibilite.

\section{Cognition simulee}
Des chaines de Markov d'ordre un apprennent les sequences de clics par site ; la satiete croit avec
les visites et decroit exponentiellement sur 10 minutes ; la pression de nouveaute favorise
les elements jamais vus ; l'anticipation du geste score les candidats par angle, distance et
habitudes ; les roles sociaux (meneur, suiveur, solitaire) et les humeurs (calme, grognon,
joueur, fier) rendent le comportement lisible.

\section{Conclusion}
Un comportement vivant n'a pas besoin d'apprentissage : des machines a etats composables et les bonnes
formules classiques --- avec des substitutions exactes et mesurees --- suffisent.
\end{document}
